\textbf{结论名称：} 将军饮马模型（两定点$+$一直线，求$PA+PB$最小值）\\
\textbf{适用条件：} 定点$A,B$在直线$l$的同侧，动点$P$在直线$l$上。\\
\textbf{变量说明：} 定点$A,B$，直线$l$，动点$P\in l$；$A'$为$A$关于$l$的对称点。\\
\textbf{核心公式：} 连接$A'B$交$l$于$P$，则$PA+PB$取得最小值，且
\[
\boxed{(PA+PB)_{\min}=A'B}
\]
即当$P$为$A'B$与$l$交点时取等号，此最小值等于$A'B$的长度。\\
\textbf{使用场景：} 中考数学中的最短路径问题，光线反射问题，解析几何中求两线段和的最小值。\\
\textbf{简单例子：} 以直线$l$为$x$轴，取$A(1,2),\,B(3,4)$。作$A$关于$x$轴的对称点$A'(1,-2)$，则
\[
A'B=\sqrt{(3-1)^2+(4+2)^2}=\sqrt{4+36}=2\sqrt{10}.
\]
由$A'(1,-2),\,B(3,4)$得直线$A'B:\,y+2=3(x-1)$，令$y=0$得交点$P\bigl(\frac{5}{3},0\bigr)$，此时$PA+PB=2\sqrt{10}$。